\subsection{选定的素槽位与行计数}\label{sec:selection}

使用方程~\eqref{eq:part-II-bootstrap} 中的 \(\Delta\) 与 \(\kappa\)：
\[
 0<\Delta\le\frac1{24},\qquad
 \kappa=2\beta_*-1=\frac34+2\Delta\in(3/4,5/6],\qquad
 \alpha:=\frac56.
\]
参数 \(\kappa\) 保持动态；\(\alpha\) 是来自六次幂放大的固定斜率。由于 \(\kappa<1\) 且
\(\beta_*=(1+\kappa)/2\)，引理~\ref{lem:plain} 的正槽位假设以等号成立。方程~\eqref{eq:part-II-bin-ceiling}
对当前的 Part~II 分箱给出 \(\delta\le\kappa\le\alpha\)，包括
可能的端点 \(\delta=\alpha\)。

固定一个动态与幅度分箱，并固定外部物理
Mellin 参数。用 \(z_{\rm phys}\) 表示固定的第三个 Mellin
参数，于是物理槽位为 \(Q_i(u;z_{\rm phys})\)。
所有 \(g_i\)，从而 \(q\)，现在都在其行和内固定。
在基 \(U\) 处，槽位 \(i\) 的长度为
\(w_i=\ell_i/d\)，可用总长度为 \(\ell/d\)。
对于请求的长度 \(0\le z\le\ell/d\)，将正的
\(g_i\) 按递减顺序排列，并按该顺序以分数方式填充 \(z\)。
所获得的指数至少为 \(qz\)：若正槽位足以
填满 \(z\)，其初始加权平均至少为所有槽位的
平均 \(q\)；否则将它们全部保留给出
增益 \(q\ell/d\ge qz\)。去掉那个可能为分数的
槽位至多损失 \((\max_i w_i)\delta/2\)。因此，长度至多为 \(z\) 的
整正槽位的一个固定子集满足
\[
 \left|\prod_{i\ {\rm selected}}Q_i\right|^2
 \ge U^{2qz-\delta\max_iw_i}.
\]
当需要一个严格矩不等式时，我们先将 \(z\)
替换为 \(z-\nu_0\)，其中 \(\nu_0>0\) 为固定的小量，或当
\(z\le\nu_0\) 时不选取任何槽位。在此显示式中所致的损失至多为
\(2q\nu_0+\delta\max_iw_i\)，除所述的零容量
邻域之外，在那里将使用未加权矩。两种损失
都可以做得比任何预先给定的正幂更小。

选定的因子恰好具有矩所要求的系数类。
如有必要，对两个见证因子（包括其
剖面）取共轭，并将其行特征写为
\(\psi(n)=\nu(n)\overline{\chi_n(u)}\)，\(\nu\in\Theta\)。
相对于此行，一个物理素数的系数为
\[
 \overline{\nu(p)}1_{p\in1_T}
 =\frac1{|T|}\sum_{\theta\in\widehat T}
       (\overline{\nu}\theta)(p).
\]
这是 \(\Theta\) 中成员的固定有限组合；
它独立于固定行和中的移动行。
没有要求 \(\eta(p)=1\)。
因子 \((q_p/P_i)^{z_{\rm phys}-1}\) 是其光滑剖面的一部分。
在掩码之前，底层的素支撑保持不相交。
带标记的矩允许此负的公共行定向。
为应用引理~\ref{lem:plain}（其行特征具有正
定向），对两个 plain 见证因子与所有选定的素因子的
整体乘积取共轭。其绝对平方
不变，且其公共行特征变为
\(\overline{\psi(n)}=\overline{\nu(n)}\chi_n(u)\)。
两个见证剖面与每个选定的槽位剖面都被共轭，
且每个选定的槽位系数变为
\(\nu(p)1_{p\in1_T}\)，再次是 \(\Theta\) 中成员的
固定有限组合。取共轭保留所有零延拓与底层
素支撑。它还保持诱导特征是否
属于 \(\Theta\)，因为 \(\Theta\) 是一个群。
见证高度与物理高度不必相等。对逐行见证参数
应用引理~\ref{lem:smooth-calculus}，
在固定物理参数与选定指标之后。
对这些参数的求导只插入对数剖面
权重。代价是 \(1+T_1\) 的固定幂，对
移动行与外部标签一致。每个保留的大行都在
\(\Theta\) 之外诱导，因为它在 \(S\) 之外的某个素数处分歧。
因此引理~\ref{lem:plain} 的 \(z>0\) 族条件成立。

对于实际的 plain 长度 \(m\le1/2\)，其在有效行宽为一处
来自引理~\ref{lem:plain} 的最大零损失
容量为
\begin{equation}
 z_P(m)=\frac{1-2m}{6\kappa}
       =\frac{1-2m}{9/2+12\Delta}.
 \label{eq:plain-zero-loss-capacity}
\end{equation}
的确，该矩被应用于 plain 见证的两份拷贝，
因此其条件为 \(2m+6\kappa z\le1\)。
对于长度为 \(r<1\) 的逆见证，令
\(z_M(r)=(1-r)/2\)。我们只在引理~\ref{lem:marked} 中的第二个
严格不等式也具有固定余量处使用此容量。

\begin{proposition}[来自见证的行计数]\label{prop:detector-counts}
设 \(\mathcal B\) 为一个固定的动态与幅度分箱，由保留的
无六次幂因子的物理行 \(q_u\asymp U=Z^d\) 组成，其中
\(a>51/100\)、\(0<\delta=2a-1\le\alpha\)，平均幅度
\(q\in[0,\delta/2]\)。令 \(x=q/\delta\)，故 \(0\le x\le1/2\)，
并定义
\[
 D_x=3-\frac{17x}{9},\qquad
 P_x=\left(2-\frac{8x}{9}\right)(1-x).
\]
假设可用的总素长度以固定的正量超过 \(7/37\)。
对每个 \(t\in[1,3/2]\) 与每个
\(\epsilon>0\)，可仅用 \(\epsilon,\Delta\) 及有界实范围选取
容量减量与槽位网格，使得
\[
 \#\mathcal B\ll_{\mathcal A,\epsilon}
 U^{\max\{R_{\rm short}(t),L(t)\}+\Delta/4+\epsilon}
 (1+T_1)^{A_{\mathcal A}},
\]
其中 \(A_{\mathcal A}<\infty\) 对移动行一致，且
\[
 R_{\rm short}(t)=1-\delta+\frac{\delta P_x}{D_x}(3/2-t),
 \qquad
 L(t)=1-\delta+(\alpha-\delta)(t-1).
\]
该估计对命题~\ref{prop:detector-witness} 的逐行见证成立，
并带有其高度条件。若未选取任何素槽位，则 \(t=1\) 转而给出
\[
 \#\mathcal B\ll_{\mathcal A,\epsilon}
 U^{1-2\delta/3+\epsilon}(1+T_1)^{A_{\mathcal A}},
\]
无需素供给假设。在零逆容量处使用
引理~\ref{lem:inverse-amplification}；在零 plain 容量处使用
引理~\ref{lem:plain} 的零槽位情形。
\end{proposition}

\begin{proof}
按见证表示与二进对进行细分。选择数
为 \(O_{\mathcal A}((\log U)^2)\)；余下的
逐行光滑参数由前述 Sobolev
论证处理。在一个细分中进行，仍将其记为
\(\mathcal B\)，并设 \(r,m\) 为命题~\ref{prop:detector-witness} 中的
实际长度。每当所选
槽位满足相应的矩假设且总请求
容量为 \(z\) 时，其尖峰与各个见证
尖峰蕴含，在扣减预备损失之后，
\[
 \begin{array}{ll}
 \#\mathcal B
 \ll U^{1-\delta r-2qz+\epsilon}(1+T_1)^{A_{\mathcal A}},
 &\text{来自引理~\ref{lem:marked}},\\[2mm]
 \#\mathcal B
 \ll U^{1-2\delta m-2qz+\epsilon}(1+T_1)^{A_{\mathcal A}},
 &\text{来自引理~\ref{lem:plain}}.
 \end{array}
\]
在第二行中使用 \(S_m\) 的两份拷贝，因此分母
为 \(|S_m|^4\)，而非 \(|S_m|^2\)。有效宽度为一，
因为扭 \(\nu\) 在行和内固定且无
移动的导子根式。每个公式都包含任意
小的容量、取整、矩与见证损失。

\paragraph{未选素数的逆见证。}
每当未选取任何逆槽位时，转而使用引理~\ref{lem:inverse-amplification} 的
六次幂放大。它适用于
这些物理行，因为它们的理想赋值至多为
五，且其原始零延拓不变。在当前
细分之后，\(U,t,D,D_*\) 是公共的，且逆基
剖面为
\[
 W_{\rm base}(y)=W_1(y)V_{\le}(Dy/D_*).
\]
其固定的对数半范数一致有界：第二个
因子的导数在其缩放自变量位于
固定紧区间处有支撑。余下的逐行参数为
固定紧区间内的 \(\sigma\) 以及纯扭
\(-(\gamma-\nu)\)，其绝对值至多为 \((3I+1)T_1\)。
因此引理~\ref{lem:inverse-amplification} 的逐行断言
在该高度范围内适用。将其因子
\((1+(3I+1)T_1)^A\) 替换为
\((3I+2)^A(1+T_1)^A\) 只改变一个固定常数。

为清晰起见，即便 \(r\) 随 \(U\) 趋于一，
此用法也是一致的。若 \(0\le r\le R_0\)，带有预备损失 \(\epsilon_0\) 的
方程~\eqref{eq:amplified-note} 的指数为
\[
 \max\left\{1,\frac{1+5(1+c)r}{6}\right\}
       +(1+r)\epsilon_0
 \le e(r)+\frac{5cR_0}{6}+(1+R_0)\epsilon_0,
 \qquad e(r)=\max\left\{1,\frac{1+5r}{6}\right\}.
\]
对于请求的损失 \(\epsilon_m>0\)，例如取
\(c=3\epsilon_m/(10\max\{R_0,1\})\) 与
\(\epsilon_0=\epsilon_m/(4(1+R_0))\)。它们在
\(U\) 之前固定，且两个额外项至多为 \(\epsilon_m/2\)。
因此除以逆见证尖峰给出
\begin{equation}\label{eq:no-slot-inverse-count}
 \begin{aligned}
 \#\mathcal B&\ll
 U^{e(r)-\delta r+\epsilon}(1+T_1)^{A_{\mathcal A}},\\
 e(r)-\delta r&=
 \begin{cases}
  1-\delta r,&0\le r\le1,\\
  1-\alpha+(\alpha-\delta)r,&r\ge1.
 \end{cases}
 \end{aligned}
\end{equation}
特别地，这不是对具有收缩余量 \(1-r\) 的严格带标记
矩的应用。

\paragraph{无需选定素数的情形。}
我们首先处理见证已在未选取素数的情况下给出
所需计数的情形。
对 \(r\ge1\)，方程~\eqref{eq:no-slot-inverse-count} 使用有效
矩指数 \((1+5r)/6=1-\alpha+\alpha r\)。
除以逆尖峰给出
\(1-\alpha+(\alpha-\delta)r\)。
由于 \(\delta\le\alpha\) 且 \(r\le t+O(\epsilon)\)，这至多
为
\begin{equation}
 L(t)=1-\alpha+(\alpha-\delta)t
     =1-\delta+(\alpha-\delta)(t-1)
 \label{eq:long-count}
\end{equation}
精确到 \(O(\epsilon)\)。这包括 \(r=1\)，此时
方程~\eqref{eq:no-slot-inverse-count} 中的两个
分支一致。
若 \(m\ge1/2\)，引理~\ref{lem:plain} 的零槽位情形
给出计数指数 \(1-2\delta m\le1-\delta\)。
由于 \(R_{\rm short}(t)\ge1-\delta\)，这也满足
所述的界。因此，对于余下的论证我们可以假设
\[
 r<1,\qquad m<1/2.
\]
此时两个容量 \(z_M(r)\) 与 \(z_P(m)\) 都是正的；
当其中任一相对于固定减量太小时，我们将在下文使用其
未加权估计。

\paragraph{比较正容量。}
对于逆容量 \(z_M(r)\)，所得到的理想指数
为下文的 \(A_I(r)\)。为比较 plain 指数，先将
\(z_P(m)\) 替换为其在 \(\Delta=0\) 处的值，即
\(2(1-2m)/9\)。由于实际的 \(m\) 至少为
\(t-r-O(\epsilon)\)，所得理想比较指数为
\begin{equation}
 \begin{aligned}
 A_I(r)&=1-\delta\{x+(1-x)r\},\\
 S_t(r)&=1-\delta\left\{\frac{4x}{9}
          +\left(2-\frac{8x}{9}\right)(t-r)\right\}.
 \end{aligned}
 \label{eq:selected-counts}
\end{equation}
实际的 plain 指数
\(1-2\delta m-2q(1-2m)/(9/2+12\Delta)\) 及其基线
版本都关于 \(m\) 递减：它们的导数分别为
\[
 -2\delta+\frac{4q}{9/2+12\Delta}<0,\qquad
 -2\delta+\frac{8q}{9}<0.
\]
因此对 \(m\) 作 \(O(\epsilon)\) 替换是合法的。

方程~\eqref{eq:selected-counts} 中的两个仿射表达式
相交于
\[
 r_*(t)=
 \frac{(2-8x/9)t-5x/9}{D_x}.
\]
这里 \(D_x\ge37/18>0\)。直接计算，
\[
 r_*(3/2)=1,\qquad
 r_*(1)=\frac{2-13x/9}{3-17x/9}\ge\frac{23}{37},
\]
且
\[
 t-r_*(t)=\frac{(1-x)t+5x/9}{D_x}.
\]
最后一个表达式关于 \(t\) 递增，在
\(t=3/2\) 处等于 \(1/2\)，而在 \(t=1\) 处为
\((1-4x/9)/(3-17x/9)\ge1/3\)。
因此，在所述范围内，
\[
 r_*(t)\ge\frac{23}{37},\qquad
 \frac13\le t-r_*(t)\le\frac12,\qquad r_*(t)\le1.
\]

当 \(r\le r_*(t)\) 时使用 plain 计数，当
\(r_*(t)\le r<1\) 时使用逆计数，固定小的
零容量邻域内除外。在逆的一侧 \(r\ge23/37\)。
将 \(z_M(r)\) 减去 \(\nu_0\) 后，
\[
 1-r-2z\ge2\nu_0>0,\qquad
 3-2r-8z=4(1-r-2z)+(2r-1)>0.
\]
第二个不等式在其正的首项之前具有至少 \(9/37\) 的
余量。这些恰是引理~\ref{lem:marked} 在行宽为一处的两个严格宽度
条件。
逆容量至多为
\((1-23/37)/2=7/37\)。在 plain 一侧，
\(m\ge1/3-O(\epsilon)\)，故其容量至多为
\(2/27+O(\epsilon)\)。在 \(d=h\) 处可用长度为
\[
 \frac{\ell}{h}=\frac8{39},\qquad
 \frac8{39}-\frac7{37}=\frac{23}{1443}>0.
\]
因此，所假设的正供给余量与足够精细的网格
允许刚刚作出的每一次选取。

在 plain 一侧，将基线容量替换为
方程~\eqref{eq:plain-zero-loss-capacity} 中的实际容量，会使
计数指数提高
\begin{equation}
 \begin{split}
 2q(1-2m)
 \left(\frac29-\frac1{9/2+12\Delta}\right)
 &=\frac{24q(1-2m)\Delta}{(9/2)(9/2+12\Delta)}\\
 &\le\frac{\Delta}{4}+O(\epsilon).
 \end{split}
 \label{eq:capacity-comparison}
\end{equation}
对于该不等式，使用 \(m\ge1/3-O(\epsilon)\)、\(2q\le1\)，
以及每个分母的下界四。\(O(\epsilon)\) 项中的
常数是一致的。

加权平均
\begin{equation}
 \begin{aligned}
 R_{\rm short}(t)
 &=\frac{(2-8x/9)A_I(r)+(1-x)S_t(r)}{D_x}\\
 &=1-\delta+\frac{\delta P_x}{D_x}(3/2-t)
 \end{aligned}
 \label{eq:short-crossing-count}
\end{equation}
独立于 \(r\)：\(r\) 的系数相消，且
权重之和为 \(D_x\)。它是在 \(r_*(t)\) 处的公共值。
由于 \(S_t\) 关于 \(r\) 递增而 \(A_I\) 递减，
所选的短计数至多为 \(R_{\rm short}(t)+\Delta/4+O(\epsilon)\)。

\paragraph{小容量与结论。}
若 \(r<1\) 但 \(z_M(r)\le\nu_0\)，使用 \(e(r)=1\) 的同一无槽位
估计给出
\[
 1-\delta r\le1-\delta+2\delta\nu_0.
\]
若 \(z_P(m)\le\nu_0\)，则
\(1-2m\le6\kappa\nu_0\le6\nu_0\)，其计数至多为
\(1-\delta+6\delta\nu_0\)。由于 \(R_{\rm short}(t)\ge1-\delta\)，
在给定的 \(\epsilon\) 内选取 \(\nu_0\) 可保持
短界。最后，\(r=1/2-O(\epsilon)\) 与 \(t\ge1\)
迫使 \(m\ge1/2-O(\epsilon)\)，故同样的未加权 plain
估计适用。不请求任何负容量。
连同无需选定素数的情形，在加上有限多个细分之后，
这就证明了所述的选定计数。

最后，不使用素槽位并取 \(t=1\)。对逆的
一侧使用方程~\eqref{eq:no-slot-inverse-count}，包括
对所有 \(r\le1\) 的 \(e(r)=1\)；对 plain 一侧使用
引理~\ref{lem:plain} 的零槽位情形。取选择增益为零，
即在两条比较式中取 \(x=0\)，作同样的短计算，
给出 \(R_{\rm short}(1)=1-2\delta/3\)；长界为
\(L(1)=1-\delta\)。两个输入都没有素供给假设，
从而证明了最后的断言。
\end{proof}

\begin{remark}[超出 Part II 分箱上限的未选定逆见证]
方程~\eqref{eq:no-slot-inverse-count} 的推导只用到
命题~\ref{prop:detector-witness} 与
引理~\ref{lem:inverse-amplification} 中的逆尖峰。因此它可以对满足这些假设的
共享探测器的任意实例重复，独立于
当前的 Part~II 分箱上限。从而，对每个固定的
\(t\in[1,3/2]\)，方程~\eqref{eq:no-slot-inverse-count} 中的界
在具有来自该命题的实际见证的固定保留行动态
分箱的每个固定表示/二进细分上成立，每当
\[
 a>51/100,\qquad 1/50<\delta=2a-1\le1,
\]
且 \(r\) 位于给定的有界非负范围内。它使用
公共逆剖面 \(W_{\rm base}\)、逐行参数界，
以及前述证明中用到的见证损失与高度假设，
但既不要求
\(\delta\le\alpha\) 也不要求素供给假设。此结论
不包含下界分箱 \(a=51/100\) 中的所有行，这些行不必
具有实际见证。
\end{remark}

命题~\ref{prop:detector-counts} 中的公式描述了下文所需的
唯一行指数。基数因子
\((1+T_1)^{A_{\mathcal A}}\) 在此保持显式。我们验证
引理~\ref{lem:late-height-closure} 所要求的
额外依目标的高度上限。设 \(\epsilon_{\rm ht}>0\) 为
已与实损失一同选取的有限多个探测器高度容差
的最小值。对固定目标，在内部剖面
阶数固定之后，设 \(A_{{\rm ht},\eta}\ge0\) 支配其有限多个
高度阶。令
\[
 \tau_{0,\eta}:=
 \frac{d_{\min}\epsilon_{\rm ht}}{20(A_{{\rm ht},\eta}+1)}>0.
\]
对 \(0<\tau\le\tau_{0,\eta}\)，\(T_1=Z^\tau\)，且 \(Z\) 充分大
时，不等式 \(U\ge Z^{d_{\min}}\) 给出
\[
 (1+T_1)^{A_{{\rm ht},\eta}}
 \le 2^{A_{{\rm ht},\eta}}Z^{d_{\min}\epsilon_{\rm ht}/20}
 \le U^{\epsilon_{\rm ht}/10}.
\]
因此此上限强制所有前述的探测器高度假设。
它在外部尾部阶与 \(Z\) 之前固定；增大该
阶只改变外部检验半范数，而不改变这些内部高度
阶。于是引理~\ref{lem:late-height-closure} 作出最终高度
选择，而不改变正实余量或槽位网格。

\section{八分之七界}\label{sec:conclusion}

我们完成 Part~II 中所假设的矛盾。回顾
方程~\eqref{eq:part-II-bootstrap} 与~\eqref{eq:part-II-bin-ceiling}
可知
\[
 0<\Delta=\beta_*-\frac78\le\frac1{24},\qquad
 \kappa=\frac34+2\Delta\le\frac56,\qquad
 \delta=2a-1\le\kappa
\]
对每个保留的分箱成立。补偿低估计是
命题~\ref{prop:probe-low}。对同一物理探针，我们先将
其高展开式的主项正规化，然后界定余下的
行，并验证命题~\ref{lem:continuation-criterion} 所要求的
与目标无关的余量。

几何量与 Mellin 指数为
\[
 h=\frac{13}{16},\qquad \ell=\frac16,\qquad
 l_x=\frac{17}{48},\qquad l_y=\frac{23}{48},\qquad
 h=1-l_x+\ell,
\]
\[
 C(s)=C_{\mathrm{II}}(s)=s-\frac{11}{16},
 \qquad C(7/8)=\frac3{16}.
\]
这些就是方程~\eqref{old-eq:5.1}
与~\eqref{eq:part-II-mellin-target} 中的值。将节~\ref{sec:local-compensation} 中
称为 \(x\) 的局部变量记作 \(s\)；下文的幅度比
\(x=q/\delta\) 是另一个不同的实数。

使用定义~\ref{def:stage-high-data}，取
\[
 \mathscr I_\eta(Z)=I_{\eta,\mathrm{modified}}(Z),\qquad
 \mathfrak H_{\eta,u,Z}(s,w,z)
     \ \text{如方程~\eqref{eq:holomorphic-selected-tuple} 中}.
\]
在共享解析引理中取 \(\sigma_0=7/8\)。
节~\ref{sec:local-compensation} 验证这些数据的每一部分：
方程~\eqref{eq:compensated-poisson-identity} 是对
方程~\eqref{eq:compensated-probe-definition} 中独立定义的有限表达式的绝对
收敛高恒等式，且全修正项在两个所需的 Euler 区域内
全纯，并满足
方程~\eqref{eq:stage-all-height-majorant}。排除集、物理
行掩码与校准均不变。特别地，在此全修正项之外没有
额外的 Euler 因子。其标量
分母为 \(L_F^S(s,\eta\overline{\chi_\bullet(u)})\)，其
Gaussian 为 \(\Phi(s+z-1)\)。

在下文的估计中，\(0<e<10^{-3}\) 是一个容许的探测器宽度。
它的最终选取，连同余下的实损失，在
结尾的选择次序中作出。

\subsection{主正规化因子}\label{sec:principal-signal}

我们验证引理~\ref{lem:principal-residue-interface} 的两个修正假设。
首先，方程~\eqref{eq:principal-local-tuple-bound} 给出
\[
 |\mathfrak H_{\eta,1,Z}(s,w,z)|\ll Z^{\ell\Re z}
\]
在第二个 Euler 区域内的每个固定实盒上对所有虚部一致成立。
特别地，它在方程~\eqref{eq:principal-correction-bound} 中的整个
矩形上以高度阶零成立。
这是对完整元组修正项的界。

对于留数值，定义
\[
 S_i(Z)=\sum_{p\in\mathcal P_i(Z)}W_i(q_p/P_i)q_p^{-5/6},
 \qquad P_i=Z^{\ell_i}.
\]
引理~\ref{lem:ray-prime-normalizer}（包括其关于删去
一个固定有限集的断言）给出
\[
 S_i(Z)\sim\frac{P_i^{1/6}}{|T|\log P_i}
       \int_0^\infty W_i(y)y^{-5/6}\,dy.
\]
每个首项常数都是正的。由于槽位数固定，
对所有充分大的 \(Z\)，所有 \(S_i(Z)\) 都是正的，允许
一个依赖于固定数据的公共下阈值。令
\begin{equation}\label{eq:principal-scalar}
 A_T(Z)=(-1)^K Z^{-\ell/6}\prod_{i=1}^K S_i(Z).
\end{equation}
利用 \(\log P_i=\ell_i\log Z\) 与 \(\sum_i\ell_i=\ell\)，我们得到
\[
 A_T(Z)\sim
 \frac{(-1)^K}{|T|^K(\log Z)^K}
 \prod_{i=1}^K\left\{\frac1{\ell_i}
       \int_0^\infty W_i(y)y^{-5/6}\,dy\right\}.
\]
从而在该范围内 \(A_T(Z)\ne0\)，且
\[
 |A_T(Z)|\asymp_{T,K,(\ell_i,W_i)}(\log Z)^{-K},
 \qquad |A_T(Z)|^{-1}\ll_\epsilon Z^\epsilon
\]
对每个 \(\epsilon>0\)。下阈值可以依赖于目标的
有限排除集；这里没有使用移动的射线导子中的
一致素数渐近。

在主第二 Euler 区域上，节~\ref{sec:local-compensation}
定义 \(\mathcal B_p=G_p/H_p\)，其中 \(H_p\ne0\)，且
方程~\eqref{eq:principal-local-approximation} 给出
\(\mathcal B_p=-1+O(q_p^{-7/8})\)，对所有虚部与
目标单位相位一致成立。将方程~\eqref{eq:principal-slot-factorization} 中的
主因子分解在 \(w=1,z=1/6\) 处应用，
即便原先对 \(P_p\) 的商未定义它仍然成立。
它给出
\[
 \mathfrak H_{\eta,1,Z}(s,1,1/6)
 =H_\eta(s)\prod_{i=1}^K
       \left(\sum_{p\in\mathcal P_i(Z)}
           W_i(q_p/P_i)q_p^{-5/6}\mathcal B_p\right),
\]
其中 \(H_\eta\) 是方程~\eqref{eq:shared-principal-data} 中
针对同一排除集 \(S\) 的函数。由 \(W_i\) 的非负性，其第 \(i\) 个
槽位等于
\[
 -S_i(Z)\{1+\rho_i(s)\},\qquad
 |\rho_i(s)|\ll P_i^{-7/8}
 \quad(\Re s=\beta_*+e,\ \Im s\in\mathbb R).
\]
的确，绝对局部误差之和至多为
\(P_i^{-7/8}S_i(Z)\) 的固定倍数。选取任意预目标数
\[
 0<\kappa_P<\frac78\min_i\ell_i.
\]
由于 \(K\) 固定，\(\mathcal R_{\eta,Z}(s):=\prod_i(1+\rho_i(s))-1\)
在该整条直线上满足 \(|\mathcal R_{\eta,Z}(s)|\ll Z^{-\kappa_P}\)。
因此精确的留数修正为
\[
 \mathfrak H_{\eta,1,Z}(s,1,1/6)
 =H_\eta(s)Z^{\ell/6}A_T(Z)\{1+\mathcal R_{\eta,Z}(s)\},
\]
即方程~\eqref{eq:principal-correction-residue}，其中
\(A_\eta=A_T\) 且 \(\mu=\kappa_P\)。此因子分解只在
主留数处断言；它不是用于一般
围道移动的修正。

设 \(c_S>0\) 为方程~\eqref{eq:shared-principal-data} 中的标量，
其正性是引理~\ref{lem:principal-residue-interface} 的一部分。
对充分大的 \(Z\)，定义正规化物理探针
\[
 J_{\mathrm{II},\eta}(Z)
   =\frac{I_{\eta,\mathrm{modified}}(Z)}{c_S A_T(Z)}.
\]
它不同于 \(J_{\mathrm I,\eta}\)，且其分子与
正规化因子都独立于分析高度 \(T_1\)。
命题~\ref{prop:probe-low} 与 \(A_T^{-1}\) 的次幂界
给出，对每个 \(\epsilon>0\)，
\begin{equation}\label{eq:part-II-normalized-low}
 |J_{\mathrm{II},\eta}(Z)|
 \ll_{\eta,\epsilon} Z^{C(7/8)+\epsilon}.
\end{equation}
同一个标量被用于高比较。
